% chapters/chap1.tex
% \include で読み込まれるファイル
% \begin{document} は書かない

\chapter{\texttt{\textbackslash include} の基本}

\texttt{\textbackslash include\{ファイル名\}} はファイルの内容を展開し、
前後に\textbf{改ページを挿入}します。
これにより、各章が必ず新しいページから始まります。

\section{\texttt{\textbackslash include} の特徴}

\begin{itemize}
  \item ドキュメント本文（\texttt{\textbackslash begin\{document\}} 以降）でのみ使用可能
  \item 前後に自動で改ページが入る
  \item ネストは不可（\texttt{\textbackslash include} したファイル内では使えない）
  \item \texttt{\textbackslash includeonly} と組み合わせて選択コンパイルができる
\end{itemize}
